% w07-why-symmetry-yaws.tex
%   7주차 §7-4. **순수 횡이동을 지시했는데도** 배가 도는 이유.
%   왼쪽 칸은 유닛 넷이 각각 내는 모멘트, 오른쪽 칸은 그것을 한 점으로 모은 것.
%   같은 것을 두 가지로 보는 그림이므로 칸을 나누고 같은 선체를 두 번 그린다
%   (사용자 지시 2026-10-17, CLAUDE.md 4-6-4).
%
%   축척 / scale
%       1 m = 24 mm. 평면도 — 화면 위가 x_b, 화면 오른쪽이 y_b.
%
%   수치의 출처 / where the numbers come from
%       W07_V_openloop.m 의 **S5** (n = 60 rad/s 넷, delta = [+45 +45 -45 -45] deg).
%       거치 부호 s = [+1 +1 -1 -1] 이므로 유닛의 힘은 s_i T_i (cos d_i, sin d_i):
%           선수 둘   (+28.205, +28.205) N
%           선미 둘   (-28.205, +28.205) N
%       네 유닛의 **횡력이 같고**(+28.205 N 넷), 종력은 선수-선미로 상쇄된다.
%       유닛별 모멘트 x_i f_iy - y_i f_ix =
%           1: 1.5( 28.205) - (-0.395)( 28.205) = +53.449
%           2: 1.5( 28.205) - ( 0.395)( 28.205) = +31.166
%           3: -1( 28.205) - (-0.395)(-28.205) = -39.346
%           4: -1( 28.205) - ( 0.395)(-28.205) = -17.064   합 +28.205 N m
%       합력의 작용점 xbar = mean(x_i) = 0.25 m = (l_f - l_a)/2,
%       그래서 N = Y xbar = 112.820 * 0.25 = 28.205 N m.
%
%   회전 방향 / the sense                        (CLAUDE.md 4-6-2)
%       N > 0 은 선수를 우현으로 — 이 평면도에서 **시계방향**이다.
%       arc(330:210:r) 은 끝 각이 작으므로 시계다.
%
%   유닛의 자리 / where the units are drawn
%       **실제 자리**다 — 선수 x = +1.5 m, 선미 x = -1.0 m. 그래서 폰툰을
%       그 자리가 들어가는 길이로 그린다 (w07-tilt-geometry.tex 와 같은 처리).
%       모멘트 팔을 눈으로 견주는 것이 이 그림의 전부이므로, 팔을 줄여 그리면
%       그림이 말하려는 것을 그림이 부정한다.

\documentclass[border=5pt]{standalone}
% gnc-style.tex — 이 볼트의 TikZ 그림이 공유하는 스타일.
%
% 저널 그림 규약을 스타일 하나로 굳혀 둔다. 그림마다 선 굵기나 화살촉을
% 다시 정하지 않는다 — 그것이 그림이 제각각으로 보이는 원인이다.
%
% 쓰는 법:  % gnc-style.tex — 이 볼트의 TikZ 그림이 공유하는 스타일.
%
% 저널 그림 규약을 스타일 하나로 굳혀 둔다. 그림마다 선 굵기나 화살촉을
% 다시 정하지 않는다 — 그것이 그림이 제각각으로 보이는 원인이다.
%
% 쓰는 법:  % gnc-style.tex — 이 볼트의 TikZ 그림이 공유하는 스타일.
%
% 저널 그림 규약을 스타일 하나로 굳혀 둔다. 그림마다 선 굵기나 화살촉을
% 다시 정하지 않는다 — 그것이 그림이 제각각으로 보이는 원인이다.
%
% 쓰는 법:  \input{gnc-style}  (figures/src/ 안에서)

\usepackage{tikz}
\usepackage{amsmath}
\usepackage{xcolor}
\usetikzlibrary{arrows.meta, positioning, calc, fit, backgrounds, decorations.pathmorphing}

% ---- 색 --------------------------------------------------------------------
% 색은 강조 하나에만 쓴다. 나머지는 검정.
\definecolor{ink}{HTML}{1A1A1A}
\definecolor{accent}{HTML}{7C3AED}
\definecolor{warn}{HTML}{B91C1C}
\definecolor{soft}{HTML}{666666}

% 기하 그림(좌표계·유도)에서만 쓰는 넷. 블록선도는 위의 흑백 + accent 로 충분하다.
\definecolor{pathc}{HTML}{0D9488}   % 계획 경로 · 항적
\definecolor{errc}{HTML}{C2410C}    % 오차 (cross-track 등)
\definecolor{alng}{HTML}{B45309}    % along-track
\definecolor{flow}{HTML}{0369A1}    % 조류 · 유속

% ---- 선과 화살촉 -----------------------------------------------------------
\tikzset{
  % 신호선. 화살촉은 작고 뾰족한 것 하나뿐이다.
  sig/.style   = {ink, line width=0.5pt, -{Stealth[length=4pt, width=3pt]}},
  wire/.style  = {ink, line width=0.5pt},
  % 강조 경로 — 파선으로 구분한다. 흑백 인쇄에서도 살아야 하기 때문이다.
  asig/.style  = {accent, line width=0.5pt, densely dashed,
                  -{Stealth[length=4pt, width=3pt]}},
  awire/.style = {accent, line width=0.5pt, densely dashed},
  % 블록. 흰 바탕, 직각, 검은 테두리.
  blk/.style   = {draw=ink, line width=0.55pt, fill=white,
                  minimum width=17mm, minimum height=8mm, inner sep=2pt, font=\small},
  % 합산점
  sum/.style   = {draw=ink, line width=0.55pt, fill=white, circle,
                  inner sep=0pt, minimum size=4.6mm, font=\scriptsize},
  asum/.style  = {draw=accent, line width=0.55pt, fill=white, circle,
                  inner sep=0pt, minimum size=4.6mm, font=\scriptsize},
  % 분기점
  dot/.style   = {ink, fill=ink, circle, inner sep=0pt, minimum size=1.5mm},
  % 신호 이름 — 선 위에 작게
  sname/.style = {font=\small, inner sep=1.5pt},
  % 그림 안의 짧은 설명. 그림 안의 글은 최소로 둔다
  note/.style  = {font=\scriptsize, soft, inner sep=1.5pt},
  anote/.style = {font=\scriptsize, accent, inner sep=1.5pt},
  % 축
  axis/.style  = {ink, line width=0.5pt},
  tick/.style  = {font=\scriptsize, soft},
}

% 캡션 한 줄 — 그림 아래 가는 선과 함께
\newcommand{\gnccaption}[2]{%
  \node[anchor=north west, text width=#1, align=left, font=\scriptsize, ink]
        at (current bounding box.south west) {\vspace{2pt}#2};}
  (figures/src/ 안에서)

\usepackage{tikz}
\usepackage{amsmath}
\usepackage{xcolor}
\usetikzlibrary{arrows.meta, positioning, calc, fit, backgrounds, decorations.pathmorphing}

% ---- 색 --------------------------------------------------------------------
% 색은 강조 하나에만 쓴다. 나머지는 검정.
\definecolor{ink}{HTML}{1A1A1A}
\definecolor{accent}{HTML}{7C3AED}
\definecolor{warn}{HTML}{B91C1C}
\definecolor{soft}{HTML}{666666}

% 기하 그림(좌표계·유도)에서만 쓰는 넷. 블록선도는 위의 흑백 + accent 로 충분하다.
\definecolor{pathc}{HTML}{0D9488}   % 계획 경로 · 항적
\definecolor{errc}{HTML}{C2410C}    % 오차 (cross-track 등)
\definecolor{alng}{HTML}{B45309}    % along-track
\definecolor{flow}{HTML}{0369A1}    % 조류 · 유속

% ---- 선과 화살촉 -----------------------------------------------------------
\tikzset{
  % 신호선. 화살촉은 작고 뾰족한 것 하나뿐이다.
  sig/.style   = {ink, line width=0.5pt, -{Stealth[length=4pt, width=3pt]}},
  wire/.style  = {ink, line width=0.5pt},
  % 강조 경로 — 파선으로 구분한다. 흑백 인쇄에서도 살아야 하기 때문이다.
  asig/.style  = {accent, line width=0.5pt, densely dashed,
                  -{Stealth[length=4pt, width=3pt]}},
  awire/.style = {accent, line width=0.5pt, densely dashed},
  % 블록. 흰 바탕, 직각, 검은 테두리.
  blk/.style   = {draw=ink, line width=0.55pt, fill=white,
                  minimum width=17mm, minimum height=8mm, inner sep=2pt, font=\small},
  % 합산점
  sum/.style   = {draw=ink, line width=0.55pt, fill=white, circle,
                  inner sep=0pt, minimum size=4.6mm, font=\scriptsize},
  asum/.style  = {draw=accent, line width=0.55pt, fill=white, circle,
                  inner sep=0pt, minimum size=4.6mm, font=\scriptsize},
  % 분기점
  dot/.style   = {ink, fill=ink, circle, inner sep=0pt, minimum size=1.5mm},
  % 신호 이름 — 선 위에 작게
  sname/.style = {font=\small, inner sep=1.5pt},
  % 그림 안의 짧은 설명. 그림 안의 글은 최소로 둔다
  note/.style  = {font=\scriptsize, soft, inner sep=1.5pt},
  anote/.style = {font=\scriptsize, accent, inner sep=1.5pt},
  % 축
  axis/.style  = {ink, line width=0.5pt},
  tick/.style  = {font=\scriptsize, soft},
}

% 캡션 한 줄 — 그림 아래 가는 선과 함께
\newcommand{\gnccaption}[2]{%
  \node[anchor=north west, text width=#1, align=left, font=\scriptsize, ink]
        at (current bounding box.south west) {\vspace{2pt}#2};}
  (figures/src/ 안에서)

\usepackage{tikz}
\usepackage{amsmath}
\usepackage{xcolor}
\usetikzlibrary{arrows.meta, positioning, calc, fit, backgrounds, decorations.pathmorphing}

% ---- 색 --------------------------------------------------------------------
% 색은 강조 하나에만 쓴다. 나머지는 검정.
\definecolor{ink}{HTML}{1A1A1A}
\definecolor{accent}{HTML}{7C3AED}
\definecolor{warn}{HTML}{B91C1C}
\definecolor{soft}{HTML}{666666}

% 기하 그림(좌표계·유도)에서만 쓰는 넷. 블록선도는 위의 흑백 + accent 로 충분하다.
\definecolor{pathc}{HTML}{0D9488}   % 계획 경로 · 항적
\definecolor{errc}{HTML}{C2410C}    % 오차 (cross-track 등)
\definecolor{alng}{HTML}{B45309}    % along-track
\definecolor{flow}{HTML}{0369A1}    % 조류 · 유속

% ---- 선과 화살촉 -----------------------------------------------------------
\tikzset{
  % 신호선. 화살촉은 작고 뾰족한 것 하나뿐이다.
  sig/.style   = {ink, line width=0.5pt, -{Stealth[length=4pt, width=3pt]}},
  wire/.style  = {ink, line width=0.5pt},
  % 강조 경로 — 파선으로 구분한다. 흑백 인쇄에서도 살아야 하기 때문이다.
  asig/.style  = {accent, line width=0.5pt, densely dashed,
                  -{Stealth[length=4pt, width=3pt]}},
  awire/.style = {accent, line width=0.5pt, densely dashed},
  % 블록. 흰 바탕, 직각, 검은 테두리.
  blk/.style   = {draw=ink, line width=0.55pt, fill=white,
                  minimum width=17mm, minimum height=8mm, inner sep=2pt, font=\small},
  % 합산점
  sum/.style   = {draw=ink, line width=0.55pt, fill=white, circle,
                  inner sep=0pt, minimum size=4.6mm, font=\scriptsize},
  asum/.style  = {draw=accent, line width=0.55pt, fill=white, circle,
                  inner sep=0pt, minimum size=4.6mm, font=\scriptsize},
  % 분기점
  dot/.style   = {ink, fill=ink, circle, inner sep=0pt, minimum size=1.5mm},
  % 신호 이름 — 선 위에 작게
  sname/.style = {font=\small, inner sep=1.5pt},
  % 그림 안의 짧은 설명. 그림 안의 글은 최소로 둔다
  note/.style  = {font=\scriptsize, soft, inner sep=1.5pt},
  anote/.style = {font=\scriptsize, accent, inner sep=1.5pt},
  % 축
  axis/.style  = {ink, line width=0.5pt},
  tick/.style  = {font=\scriptsize, soft},
}

% 캡션 한 줄 — 그림 아래 가는 선과 함께
\newcommand{\gnccaption}[2]{%
  \node[anchor=north west, text width=#1, align=left, font=\scriptsize, ink]
        at (current bounding box.south west) {\vspace{2pt}#2};}


\begin{document}
\begin{tikzpicture}[x=1mm, y=1mm, font=\scriptsize]

\def\YP{9.48}       % y_pont = 0.395 m
\def\XF{36}         % l_f = +1.5 m
\def\XA{-24}        % l_a = -1.0 m

% 선체 하나를 그리는 매크로 / one hull, drawn twice
\newcommand{\hull}{%
  \fill[ink, opacity=0.10, rounded corners=2pt] (-12.48,-30) rectangle (-6.48,42);
  \fill[ink, opacity=0.10, rounded corners=2pt] (  6.48,-30) rectangle (12.48,42);
  \draw[ink, line width=0.4pt, rounded corners=2pt] (-12.48,-30) rectangle (-6.48,42);
  \draw[ink, line width=0.4pt, rounded corners=2pt] (  6.48,-30) rectangle (12.48,42);
  \draw[soft, line width=0.35pt] (-6.48, 42) -- (6.48, 42);
  \draw[soft, line width=0.35pt] (-6.48,-30) -- (6.48,-30);
  \fill[ink] (0,0) circle (0.7);
  \node[ink, anchor=north east, inner sep=1.2pt] at (0,0) {$o_b$};
}

% 선체좌표축 / the body axes, drawn in both cells
\newcommand{\bodyaxes}{%
  \draw[ink, line width=0.4pt, -{Stealth[length=3.5pt,width=2.5pt]}]
        (0,0) -- (0,16) node[anchor=south east, inner sep=1pt] {$x_b$};
  \draw[ink, line width=0.4pt, -{Stealth[length=3.5pt,width=2.5pt]}]
        (0,0) -- (17,0) node[anchor=north west, inner sep=1pt] {$y_b$};
}

% =======================================================================
% (a) 유닛 넷이 각각 내는 모멘트
% =======================================================================
\node[ink, anchor=south, font=\scriptsize] at (2,56)
      {\textbf{(a) four equal sway forces, four unequal arms}};

\hull
\bodyaxes

%  선수 둘 — 힘 (+f_x, +f_y) 이므로 화면에서 오른쪽 위
%  / the bow pair: starboard and ahead
\foreach \p/\lab in {-\YP/1, \YP/2}{
  \node[circle, draw=ink, fill=white, line width=0.4pt,
        inner sep=0pt, minimum size=5mm] at (\p,\XF) {\lab};
  \draw[accent, line width=0.7pt, -{Stealth[length=4pt,width=3pt]}]
        (\p+1.77,\XF+1.77) -- (\p+11.31,\XF+11.31);
}
%  선미 둘 — 힘 (-f_x, +f_y) 이므로 화면에서 오른쪽 아래. 횡력은 선수와 **같다**
%  / the stern pair: starboard and astern, the same sway as the bow pair
\foreach \p/\lab in {-\YP/3, \YP/4}{
  \node[circle, draw=ink, fill=white, line width=0.4pt,
        inner sep=0pt, minimum size=5mm] at (\p,\XA) {\lab};
  \draw[accent, line width=0.7pt, -{Stealth[length=4pt,width=3pt]}]
        (\p+1.77,\XA-1.77) -- (\p+11.31,\XA-11.31);
}

\node[errc, anchor=east, inner sep=1.5pt] at (-13.5, 36) {$+53.449$};
\node[errc, anchor=west, inner sep=1.5pt] at ( 13.5, 29) {$+31.166$};
\node[errc, anchor=east, inner sep=1.5pt] at (-13.5,-24) {$-39.346$};
\node[errc, anchor=west, inner sep=1.5pt] at ( 13.5,-17) {$-17.064$};

\node[ink, anchor=north, text width=62mm, align=center, inner sep=1.5pt]
      at (2,-44) {%
  Every unit is given $n_i = 60$ rad/s, the bow pair at $\delta_i = +45^\circ$
  and the stern pair at $-45^\circ$. With $s_i$ the four \emph{sway} forces come
  out equal, $f_{iy} = 28.205$ N, and the surge forces cancel fore
  against aft.\\[3pt]
  {\color{errc}$x_i f_{iy} - y_i f_{ix}$} \,[N m] is printed beside each unit.
  The $f_{ix}$ terms cancel inside each pair, at $\pm y$, and what is
  left is\\[2pt]
  $N = f_y \sum x_i = 2\,(l_f - l_a)\,f_y$
};

% =======================================================================
% (b) 같은 힘을 한 점으로
% =======================================================================
\begin{scope}[shift={(82,0)}]

\node[ink, anchor=south, font=\scriptsize] at (2,56)
      {\textbf{(b) the same four forces, moved to one point}};

\hull
\bodyaxes

%  합력은 네 자리의 무게중심에 선다 / the resultant acts at the centroid
\fill[accent] (0,6) circle (0.9);
%  합력은 **순수 횡력**이다 — 종력이 상쇄되었으므로 화면에서 오른쪽 하나
%  / the resultant is pure sway: the surge components have already cancelled
\draw[accent, line width=1.1pt, -{Stealth[length=5pt,width=4pt]}]
      (0,6) -- (26,6);
\node[accent, anchor=west, inner sep=1.5pt, align=left] at (27.5,6)
      {$X = 0$\\ $Y = 112.820$ N};

%  xbar 치수 / the offset of that point
\draw[soft, dotted, line width=0.35pt] (0,0) -- (-19,0);
\draw[soft, dotted, line width=0.35pt] (0,6) -- (-19,6);
\draw[alng, line width=0.6pt, |-|] (-19,0) -- (-19,6);
\node[alng, anchor=east, inner sep=1.5pt, text width=22mm, align=right]
      at (-20.5,3) {$\bar{x} = 0.25$ m};

%  그 결과의 요 모멘트 — 시계방향 / the yaw it leaves, clockwise
\draw[warn, line width=0.6pt, -{Stealth[length=4pt,width=3pt]}]
      (5.20,-3.0) arc (330:210:6);
\node[warn, anchor=north, inner sep=1.5pt] at (0,-11)
      {$N = 28.205$ N m};

\node[ink, anchor=north, text width=62mm, align=center, inner sep=1.5pt]
      at (2,-44) {%
  Four equal sway forces are one sway force at the centroid of the four
  stations:
  \\[3pt]
  $\bar{x} = \tfrac14 \sum x_i = \dfrac{l_f - l_a}{2} = 0.25$ m
  \\[3pt]
  That point is not $o_b$, so a sway command given symmetrically to four
  thrusters pulls 0.25 m ahead of the origin and yaws:\\[2pt]
  $N = Y\,\bar{x} = 112.820 \times 0.25 = 28.205$ N m
};

\end{scope}

\end{tikzpicture}
\end{document}
